\documentclass[11pt]{article}
\usepackage[a4paper,margin=25mm]{geometry}
\usepackage{amsmath,amssymb,amsthm,microtype}
\usepackage[hidelinks]{hyperref}
\newcommand{\E}{\mathbb E}\newcommand{\Pp}{\mathbb P}
\newtheorem{theorem}{Theorem}\newtheorem{lemma}[theorem]{Lemma}
\title{Two contact rays replace the probability panels at large indices}
\author{}\date{29 October 2026}
\begin{document}\maketitle

The coarse acquisition need export only two numbers to the comparison at
$n\ge700$: a bound for $nD$ and a bound for the raw threshold. Its later
probability-panel calculations can be bypassed. The argument below acquires
one probability interval and a common conditional amplitude $3/5$ using
only scalar inequalities and two contact rays. It does not replace the
Fourier calculation that produced the coarse bound.

Put $\sigma=\sqrt{10}$, $p_*=(4-\sigma)/2$, and
$B_*=32.20665378$. Every terminating decimal below is an exact rational.
For a compactly supported centered variance-one law define
\[
 \rho=\E|X|^3,\quad\gamma=\E X^3,\quad\beta=\E[X|X|],\quad
 \tau(t)=\E[X\mathbf1_{X\ge t}],
\]
\[
 J=\sum_{(a,b)}\tau_{a,b}(b-a)^2,\qquad D=(3+\sigma)\rho-\gamma-3J.
\]
The sum ranges over bounded open support gaps. There are at most countably
many; all nonnegative gap sums are obtained by monotone convergence from
finite subcollections.

\begin{theorem}[Scalar post-entry acquisition]\label{a:main}
Suppose $n\ge700$ and $nD\le B_*$. Then zero is outside the closed support,
and there is an actual central gap $(-A_-,A_+)$ with positive endpoints.
For the actual sign $B=\mathbf1_{X>0}$ put
\[
 p=\Pp(B=1),\quad q=1-p,\quad v=pq,\quad m=\E X_+=\E X_-,\quad
 \ell=m/v,\quad X=\ell(B+U-p).
\]
The conditional noises $U_j$ are centered. With
\[
 \theta=\E\{X-\E(X\mid B)\}^2,\quad
 a_0=p-A_-/\ell,\quad a_1=q-A_+/\ell,\quad d=a_0+a_1,
 \quad\lambda=n\E U^2=n\theta/\ell^2,
\]
one has
\begin{gather}
 3/8<p<93/200,\qquad \theta<9/2500,\qquad d<1/250,
 \qquad \ell>499/250,\label{a:small}\\
 \rho<9/8,\qquad .066148<\beta<.2527,\qquad
 \lambda<7/20.\label{a:out}
\end{gather}
If $p\le2/5$, then additionally $\lambda<8/25$.
No contact, predecessor, exported probability interval or numerical support
radius is a premise of these conclusions.
\end{theorem}

\section{The actual gap and a short scalar localization}
Throughout this proof $D< D_0:=.047$, since $700D_0>B_*$.
If zero belonged to the support, every gap would lie on one side of it and
$(b-a)^2\le2\int_a^b|t|\,dt$. Tonelli gives $J\le\rho$, hence
$D\ge\sigma\rho-\gamma\ge\sigma-1>D_0$, a contradiction.
Centering and positive variance supply both signs and the stated gap.
Conditional means give $a_0,a_1\ge0$ and $0\le d<1$.
Write $w=1-\theta=v\ell^2$ and
\[
 Q(p)=(\sigma+1)p^2+(\sigma-1)q^2-6pq
      =(6+2\sigma)(p-p_*)^2.
\]
For the residual-noise measures
$\int f\,d\mu_\pm=\E[(X-\E(X\mid B))^2\mathbf1_{\{\pm X>0\}}f(|X|)]$,
the elementary outer-gap integral bound gives
\begin{equation}\label{a:budget}
 D\ge\ell wQ+E_0+\int k_+d\mu_++\int k_-d\mu_-,
\end{equation}
where
\[
 E_0/\ell=6w\{d(1-d)+a_0a_1\},\quad
 k_+(x)=(\sigma-1)(x+2\ell q)+6A_+,
\]
\[
 k_-(x)=(\sigma+1)(x+2\ell p)+6A_-.
\]
Indeed $J_o^+\le2\int_{A_+}^{\infty}(t-A_+)\tau(t)dt
=\E[X_+(X_+-A_+)^2]$, and similarly on the negative side. Add
$m(A_-+A_+)^2$ and substitute the conditional means. This is the common
cost identity, before discarding any positive term. It applies to endpoint
atoms and countably many gaps.

The weaker gap bound $J\le\rho+2mA_-A_+$ gives
\[
 A_-A_+\ge z:=\frac{3.162277660-1-D_0}{3}>0,\qquad w\ge A_-A_+\ge z.
\]
Here $\rho\ge1$, $\gamma\le\rho$, $m\le1/2$, and
$A_-A_+\le(\ell p)(\ell q)=w$ were used. Since all terms in
\eqref{a:budget} are nonnegative,
\[
             D\ge w^{3/2}Q/\sqrt v\ge2z^{3/2}Q.
\]
The positive-base squared tests $4z^3Q_-(.36)^2>D_0^2$ and
$4z^3Q_-(.48)^2>D_0^2$, where $Q_-$ replaces $\sigma$ by its lower
rational bound, exclude $p\le.36$ and $p\ge.48$. This uses the monotonicity
of the squared distance to $p_*\in(.36,.48)$, not an initial probability
panel.

On $.36<p<.48$, $\ell\ge2\sqrt z>1.67$ and
\[
 \frac{k_+}{\ell}\ge q\{(\sigma+5)z/w+2(\sigma-1)\}>5.2,
 \quad
 \frac{k_-}{\ell}\ge p\{(\sigma+7)z/w+2(\sigma+1)\}>5.2.
\]
The first inequality uses $x\ge A_+$ and
$A_+\ge z/(\ell p)$; the second is its signwise counterpart.
Since $w\le1$, the asserted constants follow at $q=.52,p=.36$.
The total mass of $\mu_++\mu_-$ is $\theta$, so
$D>1.67(5.2)\theta$ unless $\theta=0$. Consequently
$\theta<.006$, $w>.994$, and $\ell>1.99$.
Also $A_-+A_+\ge2\sqrt z$ gives
\[
 d\le1-2\sqrt{zv/w}<.2,
 \qquad v\ge .36(.64).
\]

Dropping $a_0a_1$ and using the lower endpoints in the kernels gives
\begin{equation}\label{a:direct}
 D\ge\ell\{wQ+6wd(1-d)+\theta\min(C_+-b_+d,C_--b_-d)\},
\end{equation}
where $C_+=3(\sigma+1)q$, $C_-=3(\sigma+3)p$,
$b_+=\sigma+5$, and $b_-=\sigma+7$.
On the current domain both noise coefficients are positive, and each
branch increases in $d$, since
$6(.994)(.6)-.006(3.162277661+7)>0$.
The inward term alone excludes $d\ge.004$:
\[
        1.99(.994)6(.004)(.996)>.047.
\]
Here $d(1-d)$ increases up to $1/2$, so no larger root was selected.

At $d=0$, the derivative with respect to $\theta$ of
$\sqrt{(1-\theta)/v}\{(1-\theta)Q+\theta C\}$ has the sign of
$(2-3\theta)C-3(1-\theta)Q$. It is positive because
$\theta<.006$, $C>6.49$, and $Q<.05$ on $[.36,.48]$.
Thus \eqref{a:direct} implies $D\ge Q/\sqrt v$.
On $[.36,3/8]$ this ratio decreases, and the squared endpoint test
$Q_-(3/8)^2>D_0^2(3/8)(5/8)$ excludes that interval.
On $[93/200,.48]$, use $Q\ge Q(93/200)>D_0/2$ and $\sqrt v\le1/2$.
This proves the probability interval in \eqref{a:small}; no incorrect
monotonicity of $Q/\sqrt v$ below $1/2$ is needed on this right interval.

Now $C\ge\min\{3(\sigma+1)(107/200),3(\sigma+3)(3/8)\}$.
The bound $D\ge\ell\theta C$, with $\ell>1.99$, gives
$\theta<.0036$ and therefore $\ell>1.996$.
The exact identity
\[
 \rho=w^{3/2}\rho_B(p)+\int(x+2\ell q)d\mu_+
                         +\int(x+2\ell p)d\mu_-,\quad
 \rho_B(p)=\frac{1-2v}{\sqrt v},
\]
and $k_+\ge(\sigma-1)(x+2\ell q)$,
$k_-\ge(\sigma-1)(x+2\ell p)$ show
\[
 \rho\le\rho_B(p)+D/(\sigma-1)
 <1.1+.047/(54/25)<9/8.
\]
Here $v\ge15/64$ and $\rho_B<1.1$ follows by a rational square test.
The exact signed-second-moment identity gives
$1-2p-2q\theta\le\beta\le1-2p+2p\theta$.
Its two endpoint evaluations give \eqref{a:out}.

\section{Noise from the same marginal budget}
The noise estimate uses the one-sided cubic inequality
$\E Z^3\le at-(t^2/a)$ for a centered $Z\le a$, $t=\E Z^2$.
It follows by expanding $\E[(a-Z)(Z+t/a)^2]\ge0$.
When $a=0$, centering forces zero noise and its term is absent.
For $\lambda_0=nq\E U_0^2$, $\lambda_1=np\E U_1^2$ and any
$\bar d\ge d$, the common budget yields
\begin{equation}\label{a:cubicbudget}
 nD\ge\ell^3\{nvQ+\lambda_0Y_0(p)+\lambda_1Y_1(p)\},
\end{equation}
\begin{align*}
 Y_0&=3(\sigma+3)p-(\sigma+7)\bar d
               +2\sqrt{6(1-\bar d)(\sigma+1)p},\\
 Y_1&=3(\sigma+1)q-(\sigma+5)\bar d
               +2\sqrt{6(1-\bar d)(\sigma-1)q}.
\end{align*}
For clarity, before arithmetic--geometric mean the extra terms are
$6nv\{d(1-d)+a_0a_1\}$ and
$(\sigma+1)\lambda_0^2/(nqa_0)+(\sigma-1)\lambda_1^2/(npa_1)$.
Use $d(1-d)+a_0a_1\ge(1-\bar d)(a_0+a_1)$ and pair each summand
with its corresponding reciprocal term. This proves
\eqref{a:cubicbudget} without an independent allocation of the conditional
variances; it is the general v13 inequality, not its old numerical table.

Take $\bar d=.004$. Both $Y_0$ and $Y_1$ increase in their respective
arguments $p$ and $q$. Exact root comparisons at $p=3/8,q=107/200$ give
$Y_0>12.8$ and $Y_1>11.8$: lower bounds $3$ and $2.6$ suffice for the two
square roots. On $p\le.4$, use $q\ge.6$ and square-root lower bound
$2.7$ to get $Y_1>12.8$. For $\lambda>0$, since $\lambda_0+\lambda_1=\lambda$,
\[
 B_*\ge nD>1.996^3(11.8)\lambda,
 \qquad \lambda<.35;
 \quad p\le.4\Longrightarrow B_*\ge nD>1.996^3(12.8)\lambda,
 \quad \lambda<.32.
\]
At zero noise the conclusions are immediate. This finishes
Theorem~\ref{a:main}.

\section{One pair of contact rays}
\begin{theorem}[A common conditional strip]\label{a:strip}
Under Theorem~\ref{a:main}, suppose also that $(X,H)$ is an attained
fixed-target pair at $C_E$ with nonnegative excess and true sharp predecessor
$C_{n-1}\le C_E$, so that the retained closed-support contact/heat bounds
below apply. Then
\[
                  -.6<U_0\le a_0,\qquad -a_1\le U_1<.6.
\]
No bound on $H$ is needed for this support conclusion.
\end{theorem}
Put $S=81/80$, $\eta=8/875$, and $c_0=819/2000<C_E$.
For $F_v(x)=\Phi(x/\sqrt v)$ the retained contact function is
\[
 \frac{H_n(H,t)}{n^{3/2}}=F_{n-1}(H-t)-F_n(H)+tF_n'(H)
                                -\frac{t^2-1}{2}F_n''(H).
\]
Closed-support contact gives
$0\le H_n(H,\pm t)+C_EQ(t;r,b)$ after discarding the nonpositive
excess term. The all-threshold heat lemma gives, for $t>\sqrt3$,
\[
 H_n(H,\pm t)\le h(t):=S^3\left\{\frac25\left(
 \frac{t^3-3t}{6}+\frac{|t^4-6t^2+3|}{72}\right)
 +\frac{t^5}{4100}+\frac{3t}{820}+\frac{17(3t^2-2)}{78720}\right\},
\]
\[
 |\partial_tH_n(H,\pm t)|<g(t):=\frac25(t^2/2+1/2+\eta).
\]
At a support contact the upper expression is
$H_n(H,\pm t)+C_EQ(t;r,b)$, where
$Q(t;r,b)=-t^3+3rt^2/2+3bt+\eta r$.
Here $r=9/8$, with $b=-.066148$ on the negative ray and $b=.2527$ on
the positive ray. These upper replacements are sign-correct.

Use just the two rational boundaries
\[
              t_-=1.996(3/8+.6),\qquad
              t_+=1.996(107/200+.6).
\]
At each, the exact checker proves the eight strict inequalities
\[
 t^2>3,\quad h+c_0Q<0,\quad Q,Q',Q''<0,
 \qquad L,L',L''<0,\quad L=g+c_0Q'.
\]
The boundary upper bounds are respectively less than $-.40$ and $-.13$.
Since $Q'''=-6$ and $L''=2/5-6c_0<0$, all continuation signs persist
along the entire ray. The derivative of the actual upper function is
strictly below $L$, because $Q'<0$ and $C_E>c_0$. Thus both exterior
rays, including boundary atoms, contain no support point. This differentiates
the exact heat expression, not its absolute-value majorant.

Finally $X>-t_-$ implies $U_0=X/\ell+p>-(3/8+.6)+p>-.6$.
Similarly $X<t_+$ implies $U_1=X/\ell+p-1<107/200+.6+p-1<.6$.
The inner endpoints and conditional centering are the actual ones already
defined. This proves Theorem~\ref{a:strip}.

\section*{The resulting short chain, and its limits}
The selected coarse-entry theorem exports $nD\le B_*$ and
$|H|\le1.83605680<1.84$. The two theorems above then supply, on one
interval $(.375,.465)$, the support and noise premises for the v15 local
unions: use the left union when $p\le.4$ and the right union otherwise.
Their strict signed conclusions exclude an acquired nonnegative witness
at every $n\ge700$.

No v10 numerical localization, one-pass panel export, v11 left-strip test,
v13 degree-eight noise table or low/high-probability exclusion is needed
on this route. The underlying common-cost and one-sided-cubic lemmas remain
and have been displayed here. The manuscript's two coarse passes, four
signed passes and complete Fourier inputs remain ancestors of $B_*$ and
the raw-threshold ceiling. Their internal panels and intermediate support
gaps cannot be erased merely because these outputs are enough downstream.
The strict predecessor, attainment, reflection, and finite base are still
separate obligations. In particular this is not an unconditional cutoff at 700.

The companion audit gives source locations and a dependency graph. Its
arithmetic script uses exact rationals and exact $\mathbb Q(\sqrt{10})$
actual-law checks. A rare conditional atom at $-.8$ satisfies the geometric
budget and even the old radius $3.14$, but violates the $.6$ strip; it
explicitly shows why contact cannot be dropped.
\end{document}
